\section{Related Work}
\label{sec:related}

\textbf{Excavator modelling.} Mechanical and hydraulic models of the excavator manipulator are
reviewed in \cite{xu_review_2016}. Multibody formulations of the front attachment
\cite{noauthor_dynamical_2017} and dynamic simulations that determine joint reaction forces and
cylinder driving forces \cite{dang_tan_dynamic_2020} describe the mechanical side; nonlinear
models of the hydraulic system are simulated and validated in \cite{casoli_simulation_2015}.
Estimating the bucket-tip position from the measured boom, stick and bucket motion is studied in
\cite{sun_analysis_2020}. This work uses a single-joint boom model with parameters from a
Simscape multibody/hydraulic model, which keeps reinforcement-learning training fast.

\textbf{Energy efficiency.} Energy saving in excavators is mostly addressed in the hydraulic
circuit, for example with hydraulic hybrid architectures \cite{casoli_hydraulic_2016} or with
energy recovery that combines a flywheel and flow regeneration on the boom
\cite{li_novel_2020}. This paper instead asks how much energy the \emph{controller} of an
unmodified boom can save, and measures it against a computed energy floor.

\textbf{Learning-based excavator control.} Reinforcement learning has been combined with
multibody dynamics to automate excavation \cite{kurinov2020excavator}, and applied to the
control of hydraulic excavator arms \cite{egli2022excavator}. This work uses Soft Actor-Critic
\cite{haarnoja2018sac}, residual reinforcement learning on top of a conventional controller
\cite{johannink2019residual}, and dynamics randomisation for robustness to model error
\cite{peng2018dr}, and compares the learned policies with tuned PID baselines on the same test
conditions, including a plant with real machine parameters.
% NOTE: sentences describe the cited papers at title level; read each abstract and sharpen.
% noauthor_dynamical_2017 has no author in Zotero: fix it there and re-export.
